\documentclass[11pt,a4paper]{article}
\usepackage{fontspec}
\setmonofont{Latin Modern Mono}[Ligatures=TeXOff]  % -- stays two hyphens in \texttt
\usepackage[margin=20mm]{geometry}
\usepackage{microtype,tikz,booktabs,enumitem,fancyvrb}
\usepackage{alfdoc}
\usetikzlibrary{fit,positioning}
\usepackage{xurl}
\usepackage[colorlinks=true]{hyperref}
\pagestyle{empty}
\setlist{nosep,itemsep=1pt}
\newcommand{\n}[1]{\textsf{#1}}
\newcommand{\home}{\textasciitilde}
\title{\n{para}: the plan}
\author{Alfredo Canziani}
\date{14 Nov 2026}
\begin{document}

\noindent{\LARGE\bfseries \csname @title\endcsname}\\[1mm]
\csname @author\endcsname, \csname @date\endcsname\par\medskip

\noindent \n{para} is one terminal app that shows all my work and what waits on me. It lists everything I work on, with the time since I last touched it and its next step, and it keeps the queue of what waits on me now. A few items are pinned to the top, whatever else happens, and the book comes first. % Review: “I am not familiar with PMI, PMBOK guide. … Maybe you can add one these two references as well in the document, in the current version, current revision.”
The plan grows by stepwise refinement~\cite{wirth}, which project management calls progressive elaboration~\cite{pmbok}: each revision is one reading of mine, and my notes decide which part gets more detail next. This plan fixes the shape, and a design document comes next with how each part works.

\section{Why}

I have 44 notes in the vault's \n{Projects} folder, 27 project folders in \n{\home/PARA}, and 37 repos, and nothing shows them together. Each document records its rounds in its revision table, but nothing records how the work itself moves, so I cannot tell where anything stands. Today 25 PDFs wait in \n{Boox/Review} in no order, and \n{p} finds only the work that has a folder. This week I spent no time on the book.

The same holds for the small things. \n{Reminders} holds 14 open reminders, six of them late, and a snoozed reminder drops out of my sight.
% Review: “these were for you to create not to add to the report”
 My inbox was in the same state until the side pane gave it a queue: items arrive, wait for me, get handled, and close. \n{para} applies that queue to the rest of my work.

\section{The names}

Tiago Forte's PARA method~\cite{para} sorts everything into projects, areas, resources, and archives. Forte's project has an outcome and a deadline, and it goes to the archive once done. % Review: “project / that's why I initially we had a split in projects and tech projects because I called my stuff project and I don't want to call thiago stuff project so his stuff is going to be called something else …”
% Review: “I don't like this word with our other options / I think these are commitments”
What I call a project is different: I build it, it has files and usually a repo, and it runs for as long as I keep working on it. The word project stays with mine, as it did when the vault split projects from tech projects, and Forte's kind gets a name of its own:
\begin{itemize}
  \item A \textbf{project} is something I make, with files: the book, \n{AlfBOT}, the radio. It has no end date.
  \item A \textbf{commitment} is something I have to get done once: a trip, a court hearing, a visa. It ends, and then it goes to the archive. The other names we weighed were errand, task, chore, and obligation.
  \item An \textbf{area} and a \textbf{resource} keep Forte's meaning. Projects and commitments both belong to an area.
\end{itemize}
A project moves in two ways. A \textbf{session} is one sitting at the computer, with \n{Claude Code} or \n{Codex}, that changes the project, and the project note's log records it. A \textbf{round} belongs to a document: Claude writes it, I read it on the Boox and dictate notes, and Claude takes them in. The revision table records each round. Today the documents explain the project to me, and its next step comes from the sessions alone.
The \textbf{list} is every project and commitment on one screen, one row each. The \textbf{queue} is what waits on me now: documents to read, reminders, and replies.

\section{The picture}

On the left are the four places where my work already lives. \n{para} reads them and shows the list, one item, or the queue. I ask it in the terminal, and \n{AlfBOT} asks it for me. The next step that waits on me goes into \n{Reminders}, so my phone tells me. Each session writes back to the note, and each round writes back a revision. The last arrow is a proposal: a round that closes adds one line to the project's log, so what I learnt from a document moves the project.

\begin{center}
\begin{tikzpicture}[font=\sffamily\small,
  box/.style={draw,very thick,rounded corners=3pt,align=center,minimum width=34mm,minimum height=8mm},
  lab/.style={font=\sffamily\scriptsize,text=black!55}]
  \node[lab] at (0,1.0) {where my work lives};
  \node[box,draw=blue]   (v) at (0,0.2)  {vault notes, logs};
  \node[box,draw=blue]   (g) at (0,-0.6) {repos};
  \node[box,draw=blue]   (b) at (0,-1.4) {documents, Boox};
  \node[box,draw=blue]   (r) at (0,-2.2) {Reminders};
  \node[lab] at (6,1.0) {the app};
  \node[box,draw=violet,minimum height=30mm,minimum width=36mm] (p) at (6,-1.0) {};
  \node[font=\sffamily\bfseries] at (6,0.2) {para};
  \node[box,draw=violet,minimum width=26mm,minimum height=6mm] at (6,-0.6) {list};
  \node[box,draw=violet,minimum width=26mm,minimum height=6mm] at (6,-1.35) {one item};
  \node[box,draw=violet,minimum width=26mm,minimum height=6mm] at (6,-2.1) {queue};
  \node[lab] at (12,1.0) {who asks};
  \node[box,draw=aqua] (me)  at (12,-0.4) {me, in the terminal};
  \node[box,draw=aqua] (bot) at (12,-1.6) {AlfBOT};
  \foreach \s in {v,g,b} \draw[->] (\s.east) -- (p.west |- \s.east);
  \draw[<->] (r.east) -- (p.west |- r.east);
  \draw[->] (p.east |- me) -- (me);
  \draw[->] (p.east |- bot) -- (bot);
  \draw[->] (me.east) -- ++(0.4,0) |- node[lab,below,pos=.75] {a session writes to the note; a round writes a revision, then a log line} (-2.3,-3.1) |- (v.west);
\end{tikzpicture}
\end{center}

\section{A first look}

Like \n{spark}, \n{para} is one command with subcommands. \texttt{para list} prints the list, and \texttt{para show bot} opens one item by its key, which stays the same wherever the row falls. Times are relative, such as 2 days ago. The first line counts everything, so the size of my work shows in the terminal at a glance. The last column shows the past eight weeks, one mark per week, \texttt{\#} for a week with a session and \texttt{.} for one without, as GitHub's contribution grid and Anki's last review do.

\begin{Verbatim}[fontsize=\footnotesize]
$ para list
  31 projects  13 commitments  25 to read  14 reminders, 6 late

  key   item                  touched        next                     8 weeks
* bk    book                  8 days ago     next chapter             ###.##..
* bot   alfbot                2 days ago     read alfbot-plan  < me   ..######
  lwc   learning-with-claude  1 hour ago     radio/ lesson            .#######
  dmv   red-light-tickets     3 days ago     hearing date      < me   ......##
\end{Verbatim}

\section{What others do}

David Allen's \textit{Getting Things Done}~\cite{gtd} gives every project a next action, which my notes already carry as a \textbf{Next action} line, and a weekly review that finds the projects without one, which is what \n{para} builds. Allen and Forte describe methods. Others have built tools, and three of them come close to \n{para}, each with one thing for me to look at.
\begin{itemize}
  \item \n{Taskwarrior} keeps tasks in the terminal, each tagged with a project, and prints reports such as \texttt{projects} and \texttt{next}, which sorts open tasks by urgency~\cite{tw}. To look at: the \href{https://taskwarrior.org/docs/30second/}{30-second tutorial}, which shows the commands and their output.
  \item \n{Dataview} is an \n{Obsidian} plugin that turns notes into tables: each note's properties become columns, and the last change is one of them~\cite{dv}. To look at: \href{https://publish.obsidian.md/hub/04+-+Guides\%2C+Workflows\%2C+\%26+Courses/Community+Talks/YT+-+An+Introduction+to+Dataview}{An introduction to Dataview}, the first Obsidian community talk, on video.
  \item \n{cerebro} builds a terminal dashboard from git activity, agent sessions, and a status file per project, close to what \n{para} wants~\cite{cb}. To look at: \href{https://ryanparsley.codeberg.page/cerebro/print.html}{its manual}, one page.
\end{itemize}
These three go into my queue as the first items to look at.

\begin{thebibliography}{PMBOK}\raggedright
\bibitem[PARA]{para} Tiago Forte, \textit{The PARA Method: Simplify, Organize, and Master Your Digital Life}, Atria, 2023. \url{https://fortelabs.com/blog/para/}
\bibitem[GTD]{gtd} David Allen, \textit{Getting Things Done: The Art of Stress-Free Productivity}, revised edition, Penguin, 2015. \url{https://gettingthingsdone.com/}
\bibitem[Wirth]{wirth} Niklaus Wirth, \textit{Program development by stepwise refinement}, Communications of the ACM 14(4), 221--227, 1971. \url{https://doi.org/10.1145/362575.362577}
\bibitem[PMBOK]{pmbok} Project Management Institute, \textit{A Guide to the Project Management Body of Knowledge (PMBOK Guide)}, 7th edition, 2021. \url{https://www.pmi.org/standards/pmbok}
\bibitem[TW]{tw} Taskwarrior, \textit{30-second tutorial}. \url{https://taskwarrior.org/docs/30second/}
\bibitem[DV]{dv} SkepticMystic, \textit{An introduction to Dataview}, Obsidian community talk. \url{https://publish.obsidian.md/hub/04+-+Guides\%2C+Workflows\%2C+\%26+Courses/Community+Talks/YT+-+An+Introduction+to+Dataview}
\bibitem[CB]{cb} Ryan Parsley, \textit{cerebro} manual. \url{https://ryanparsley.codeberg.page/cerebro/print.html}
\end{thebibliography}

\revisiontable
\end{document}
